% =============================================================================
%  front/translator_note.tex —— 译者说明
% =============================================================================
\section*{译者说明}
\addcontentsline{toc}{section}{译者说明}

\textbf{底本。}本译本据 Ben Webster 与 Philsang Yoo 的综述《三维镜像对称》
（\textit{3-dimensional mirror symmetry}，arXiv:2308.06191v1 [math.RT]，2023 年 8 月 11 日，23 页）译出。
该文是应 \textit{Notices of the AMS} 之邀写的介绍性综述，arXiv 版比刊出版略长；
凡两版有出入处，译本一律以 arXiv 版为准。

\textbf{为什么译这篇。}三维镜像对称是过去十年「数学与三维量子场论」交叉里最核心的一条线索：
它把辛消解、辛奇点、箭图簇、超环面簇这些原本各写各的理论，解释成同一个三维 $\mathcal N=4$ 理论的
希格斯分支与库仑分支之间的对偶。对本库而言，它正好落在「镜像对称—辛几何—几何 Langlands」这条主线的交叉口上，
篇幅又短，适合作为进入这一方向的第一个中文读本。

\textbf{体例。}译文保留原文的全部编号：定理、命题、定义、例、注、问题、猜想各自按节连续编号，
与原文逐条对应，便于与原文对读。橙色的「译注」框是译者所加，只在三种情形出现：
术语背景、原文省略、原文笔误；每条都标明是哪一种。正文中的脚注是原文的，全部译出。

\textbf{关于原文第 2 页。}原文 PDF 的第 2 页是作者用 \texttt{fixme} 包留下的写作批注
（\texttt{\textbackslash listoffixmes} 的输出，共 8 条，例如「此处补一条文献」）。那是投稿时的内部备忘，
不是正文，且其中数条在定稿时已处理。本译本不排印这一页，仅在此说明。

\textbf{术语。}关键术语的译法见书末「术语表」；取舍原则是：已有通行译名的一律沿用
（symplectic resolution 辛消解、Higgs branch 希格斯分支、Coulomb branch 库仑分支），
人名与定理名不译（Springer、Koszul、Slodowy、Nakajima、Langlands），
首次出现的术语在括号里给出英文。

\textbf{许可。}原文版权归作者与 AMS。本译本是非商业的研究用中文译注本，
不代原作者授权、不另设限制；权利人通知即撤。译文、术语表与译注按本库
\texttt{COMPLIANCE.md} 处理；引用请同时给出原文与译本。

\vspace{0.5em}
\begin{flushright}
数译·SY01-30\qquad 2026 年 11 月
\end{flushright}

\clearpage
